\section{Inference}
\label{spec:inf}

\subsection{Redraw p-values}
\label{spec:inf.pval}

Fix the observed units, their order, $\sgn_u$, the logged probability tables, the weights $w_u$ and $\sigma$. A \emph{redraw} $b = 1, \dots, B$ replaces the mode of every unit of the tested family by a fresh draw from its logged table (cumulative inversion of a new uniform); the modes of the other families are kept. It gives $z_b = S_b/\sigma$ with $S_b = \sum_u v_u^{(b)} w_u$. The generator for family index $i$ and split $s$ is seeded by $[\text{seed}, 6616, 3, i, s]$.

\begin{definition}[Redraw p-values]
With $B = 9999$ and $\epsilon = 10^{-9}(1 + \abs{z})$,
\begin{equation*}
p = \frac{1 + \#\{b : \abs{z_b} \ge \abs{z} - \epsilon\}}{B + 1}, \qquad
p_d = \frac{1 + \#\{b : d\, o\, z_b \ge d\, o\, z - \epsilon\}}{B + 1},
\end{equation*}
where $d \in \{-1, +1\}$ is the direction (Section~\ref{spec:fdr.directions}) and $o$ the orientation (Section~\ref{spec:pmrt.weight}); $p_d$ is formed only where $o$ is set.
\end{definition}

\noindent The smallest attainable value is $1/(B+1) = 10^{-4}$. The p-value counts the observed statistic among $B+1$ values and is valid for every $B$ when these are exchangeable (Section~\ref{spec:thm.exact}).\claim{DEF-PMRT-PVAL}

\coderef{oran_causal.core.pmrt.engine:crt}

\subsection{Sequential stopping}
\label{spec:inf.bc}

The controlled environments use the Besag--Clifford rule. Redraws are generated one at a time up to $B = 9999$; with $h = 20$, stop at the first $k$ at which the number of exceedances reaches $h$ and set $p = h/k$ (two-sided, on $\abs{z}$). If the loop ends with a count below $h$, $p = (1 + \text{count})/(B + 1)$. The system studies use the fixed-$B$ rule above.

\subsection{Comparison tolerance}
\label{spec:inf.tolerance}

The count uses $\abs{z_b} \ge \abs{z} - 10^{-9}(1 + \abs{z})$ (and the one-sided analogue), so values that differ only by rounding count as ties. The tolerance can only increase the count and hence $p$; the test is never less conservative than the exact tie rule.
